\section{Materials and Methods}

\subsection{Problem Formulation and Observation Contract}

Let
\[
X=\{M,H,V_{\mathrm{rms}},W\}
\]
denote the available multimodal observations, where $M$ is the post-stack migrated image, $H$ is an interpreted-horizon map, $V_{\mathrm{rms}}$ is the RMS-velocity field, and $W$ is a sparse well-velocity map. The target is the depth-domain velocity model
\[
V\in\mathbb{R}^{1\times 70\times 70}.
\]
RMS velocity and the migrated image are represented on a $1000\times70$ time--lateral grid, whereas horizons, wells, and the target use the $70\times70$ depth grid. The model receives only the observation set at inference; target-derived decomposition variables are restricted to training supervision or diagnostic analysis.

\begin{table}[t]
\centering
\caption{Multimodal task definition used in the Remote Sensing evaluation.}
\label{tab:task-contract}
\begin{tabular}{llll}
\toprule
Quantity & Domain & Shape & Role \\
\midrule
Target velocity $V$ & depth & $1\times70\times70$ & prediction target \\
RMS velocity $V_{\mathrm{rms}}$ & time & $1\times1000\times70$ & kinematic/scale cue \\
Migrated image $M$ & time & $1\times1000\times70$ & structural cue \\
Horizon map $H$ & depth & $1\times70\times70$ & interface cue \\
Sparse well map $W$ & depth & $1\times70\times70$ & local velocity anchor \\
\bottomrule
\end{tabular}
\end{table}


The benchmark is constructed from the eight OpenFWI subsets FlatVel-A/B, CurveVel-A/B, FlatFault-A/B, and CurveFault-A/B \citep{deng2022openfwi}. For each reference velocity model, synthetic processed/interpreted observations are generated offline to represent a post-stack model-building workflow. This construction defines a controlled multimodal benchmark rather than a claim that these interpreted products are independent field measurements.

\paragraph{RMS velocity.}
For each lateral position, depth samples are mapped to two-way time through
\begin{equation}
t(z,x)=2\int_0^z \frac{1}{V(z',x)}\,\mathrm{d}z',
\end{equation}
and RMS velocity is computed after interpolation onto a uniform time grid:
\begin{equation}
V_{\mathrm{rms}}^2(t,x)=\frac{1}{t}\int_0^t v^2(\tau,x)\,\mathrm{d}\tau.
\end{equation}

\paragraph{Post-stack migrated image.}
OpenFWI shot records are processed through a common-midpoint workflow, normal-moveout correction using the associated RMS velocity, stacking, and post-stack time migration. The resulting migrated section provides a dense structural observation on the time--lateral grid.

\paragraph{Horizons and wells.}
The benchmark horizon map is constructed from strong vertical velocity gradients in the reference model using a fixed quantile-based rule. Sparse wells are produced by retaining selected velocity columns together with a binary observation mask. Training wells may vary across epochs, whereas validation and test wells are deterministic under the fixed well seed.

\subsection{Role-Aware Multi-Source Conditioning}

The heterogeneous inputs differ in sampling domain and information density, but the method does not impose a hard modality partition. Each modality is encoded by a modality-specific subencoder and mapped into two learned role representations,
\begin{equation}
\mathcal E_m(x_m)=\left(Z_m^{\mathrm{bg}},Z_m^{\mathrm{str}},r_m\right),
\end{equation}
where $Z_m^{\mathrm{bg}}$ and $Z_m^{\mathrm{str}}$ are background-oriented and structure-oriented feature maps and $r_m$ contains learned reliability scores. Availability and quality information are combined with these scores to produce normalized role-specific fusion weights. The fused conditions are
\begin{equation}
C_k=\mathcal G_k\left(\sum_{m\in\Omega} w_m^k Z_m^k\right),
\qquad k\in\{\mathrm{bg},\mathrm{str}\},
\end{equation}
where $\Omega$ is the set of observed modalities. $\mathcal G_{\mathrm{bg}}$ is identity in the current implementation, while $\mathcal G_{\mathrm{str}}$ is a lightweight spatial adapter.

Training-only multimodal alignment and role-calibration losses encourage the two interfaces to preserve complementary information. These auxiliary terms support representation learning but are not inference inputs and are not used as evidence for a universal decomposition of geophysical information. Detailed alignment and regularization terms can be reported in supplementary material.

\subsection{Deterministic Background Prediction}

A background network predicts
\begin{equation}
\hat B=f_{\mathrm{bg}}(C_{\mathrm{bg}})
\end{equation}
on the target velocity grid. A fixed low-pass operator $P_L$ defines the background supervision target. The background loss is
\begin{equation}
\mathcal L_{\mathrm{bg}}
=
\|\hat B-P_L(V)\|_1
+\lambda_{\mathrm{bg},2}
\|\hat B-P_L(V)\|_2^2.
\end{equation}
The purpose of this branch is operational: estimate the smooth component that can be predicted directly from the multimodal condition and provide a concrete reference for the residual task.

\subsection{Prediction-Consistent Residual Flow Matching}

The residual target is defined relative to the actual predicted background,
\begin{equation}
R_{\hat B}=V-\hat B,
\label{eq:rs_pred_residual}
\end{equation}
rather than relative to an oracle background unavailable at inference. During residual optimization, a stop-gradient copy $\bar B=\mathrm{sg}[\hat B]$ is used so the residual objective does not redefine the exclusive optimization responsibility of the background branch.

The residual vector field receives the structural condition and an embedding of the same predicted background,
\begin{equation}
c=[C_{\mathrm{str}};\psi_{\mathrm{bg}}(\bar B)].
\end{equation}
For $t\sim\mathcal U(0,1)$ and Gaussian base noise $\xi\sim\mathcal N(0,I)$, the straight interpolation path is
\begin{equation}
R_t=(1-t)\xi+tR_{\hat B},
\qquad
u_t=R_{\hat B}-\xi.
\end{equation}
The conditional Flow Matching objective is
\begin{equation}
\mathcal L_{\mathrm{FM}}
=
\mathbb E\left[
\|v_\theta(R_t,t,c)-u_t\|_2^2
\right].
\end{equation}
The current model additionally uses reconstruction-oriented terms on the one-step endpoint estimate to align the generated residual with the final composed velocity model. These terms supervise the actual residual target rather than forcing its low-pass component to vanish.

The defining consistency property follows directly from the composition
\begin{equation}
\hat V=\hat B+\hat R.
\end{equation}
Combining this with Eq.~\eqref{eq:rs_pred_residual} gives
\begin{equation}
\hat V-V=\hat R-R_{\hat B}.
\label{eq:rs_composition_identity}
\end{equation}
Thus, when training and inference use the same background reference, final reconstruction error is exactly the residual estimation error relative to that reference. This identity is the central reason for reusing $\hat B$ across supervision, conditioning, and composition; no claim about universally reduced conditional entropy is required.

\subsection{Joint Training and Observed-Only Inference}

The joint objective can be written as
\begin{equation}
\mathcal L=
\mathcal L_R+\mathcal L_{\mathrm{bg}}+\lambda_C\mathcal L_C,
\end{equation}
where $\mathcal L_C$ trains the role-aware condition interfaces, $\mathcal L_{\mathrm{bg}}$ trains deterministic background prediction, and $\mathcal L_R$ trains prediction-relative residual transport. Target-derived decomposition quantities and training-only anchors are never exposed to the public inference boundary.

At inference, the model receives only the four observed modalities and their masks. It forms $(C_{\mathrm{bg}},C_{\mathrm{str}})$, predicts $\hat B$, integrates the learned residual vector field from the Gaussian base, and outputs $\hat V=\hat B+\hat R$. The current release configuration uses 50 Euler steps for residual integration.

\subsection{Experimental Protocol}

The complete multimodal corpus contains 336,000 aligned records over eight OpenFWI subsets. A deterministic global split with fractions 0.7/0.2/0.1 and split seed 42 yields 33,600 held-out records. Test-time wells are deterministic under well seed 1234. The common held-out identity is represented by \texttt{dataset\_id}, \texttt{dataset\_name}, and \texttt{source\_sample\_index}.

RS-E01-Lite re-evaluates retained checkpoints on the same held-out membership, normalization, and evaluator. The default journal comparison contains adapted InversionNet, adapted VelocityGAN, adapted Latent U-Net, and \method. Adapted UPFWI is included only if the final comparison remains scientifically clear. These are matched-input estimator adaptations and are not presented as reproductions of the original source-native observation protocols.

The mandatory metrics are mean absolute error (MAE), root-mean-square error (RMSE), and structural similarity (SSIM) \citep{wang2004ssim}. Parameter count is reported as a supporting measure of model compactness. Low- and high-pass MAE are secondary and will only be retained if the unified evaluation confirms that their definitions and values are internally consistent across every method.
